\documentclass{article}
\usepackage[T1]{fontenc}
\usepackage{lmodern}
\usepackage{graphicx}
\usepackage{amsmath,amssymb}
\usepackage[a4paper,margin=2cm]{geometry}
\usepackage[absolute,overlay]{textpos}
\setlength{\TPHorizModule}{1mm}
\setlength{\TPVertModule}{1mm}
\usepackage[indonesian]{babel}
\usepackage{hyperref}
\setlength{\emergencystretch}{2em}
% unit: Halaman 23

\begin{document}

% === Halaman 23 (python/native) ===

23

Andaikan kriptanalis mengetahui ukuran blok = 2 karakter (16 bit), 

spasi diabaikan. 


Blok-blok cipherteks:


C1, C2, C3, C4, C5, C6, C7, C8, C9, C10, C11, C12,  C13, C14, C15, C16

Misalkan kriptanalis mengethaui posisi pesan yang berkaitan dengan 

nominal uang, yaitu blok C8 dan C9. 

Kriptanalis membuang blok cipherteks C8 dan C9 : 


C1, C2, C3, C4, C5, C6, C7, C10, C11, C12, C13, C14, C15, C16

\end{document}
